

The objective of this research is to design the architecture of and to implement a cohesive, automated compilation pipeline capable of translating high-level algorithms into low-level optimized silicon instructions. This section details the system architecture, focusing on the transformation mechanisms and the underlying platform dependencies. 

The compiler pipeline is implemented as a continuous, directed data flow that steps from the original MATLAB file into the RISC-V binary. The system architecture can be visualized through a structured sequence of state transitions and module interactions in Figure 3.1. The pipeline starts by taking in raw .m files authored by clinical researchers. These files contain standard MATLAB syntax, incorporating dynamic array allocations, calls to proprietary signal processing toolboxes, and high-level, implicitly parallel matrix arithmetic.

\begin{figure}[h]
    \centering
    \includegraphics[width=\textwidth]{figs/pipeline_architecture2.png}
    \caption{System architecture of the MATLAB-to-RISC-V compilation pipeline. Raw \texttt{.m} files are parsed and type-inferred by the ts-traversal engine. Afterwards, the MATLAB operations are mapped to RVV intrinsics via the RISCV-Matrix library. The resulting main.c file is compiled by the RISC-V GCC bare metal toolchain into a single binary, validated on the Spike ISA simulator.}
    \label{fig:pipeline-architecture}
\end{figure}


These raw files are passed into ts-traversal \cite{ts}, which is the first component of the translation pipeline. ts-traversal executes a multi-phase transformation. First, it uses the TreeSitter and TreeSitter-Matlab parsing libraries to create a concrete syntax tree (CST). This representation captures the complete syntactic structure of the input MATLAB file and the algorithms it contains. Then, the engine performs type inference on the dynamically typed input MATLAB code. The inference algorithms in ts-traversal traverse the CST to understand the maximal dimensions, memory allocation requirements, and fundamental data types of all MATLAB variables, mapping them to statically typed C constructs. The output of this phase is standard C code. The semantic intent of the MATLAB script is entirely preserved and the available vectorized and scalar functions are mapped via RISCV-Matrix library's \texttt{matrix.h} header, which contains contain highly optimized, inline RVV C intrinsics as well as reference scalar implementations of functions. This transformation effectively converts the MATLAB into heavily accelerated C code explicitly built for the parallel execution capabilities of HALO's microcontroller.

The final phase involves compiling and deploying the codebase. The generated main.c file is compiled via the RISC-V GNU Compiler toolchain. The compiler statically links the optimized vector library to produce a highly efficient RISC-V binary. Thus, the pipeline generates an executable containing the exact algorithms authored by neuroscientists in Matlab while meeting the power and latency requirements of the implanted chip. 

The validation, execution modeling, and accurate instruction-level benchmarking of the generated binaries are conducted using the Spike ISA simulator, since the actual HALO silicon is not readily available. Spike provides a fully functional, instruction-accurate model of a RISC-V core with the RVV extension, enabling precise tracking of executed instructions, vector register utilization, and overall execution cycles. This enables the continuous, iterative refinement of both the source-to-source compiler and the vectorization library without requiring constant access to the physical test hardware.

